\section{Multiprocessors}
\label{sec:mp}
\sources{Group presentation ``Multi-Processor 1''
(\nolinkurl{06-multiprocessor-part1-group.pdf}: systems, Flynn, synchronization,
consistency, superpipelining, superscalar, multithreading); Group~4
``Multiprocessor Part 2'' (\nolinkurl{06-multiprocessor-part2-group4.pdf}:
dependences, issue policies, VLIW, cache coherence); H\&P 5th ed.\ \S1.2, Ch.~3
(\S3.1--3.2, 3.4, 3.6--3.7, 3.12), Ch.~5 (\S5.1--5.6), App.~I (\S I.3).}

% =============================================================================
\subsection{Multiprocessor Systems}
\label{sec:mp-systems}

\begin{definition}[Multiprocessor]
A \textbf{multiprocessor} is a computer with two or more processors that
cooperate on tasks and share memory or have their own local memories. Related
terms: \textbf{multithreading} (several threads of execution in one program or
core), \textbf{multicore} (several processors, \emph{cores}, on one chip). The
course focuses on shared-memory, tightly coupled MIMD
multiprocessors.\HP{\S5.1}{344}
\end{definition}

\paragraph{Why.} Single processors hit three walls: the \textbf{power wall} (higher
clock rates waste energy as heat), the \textbf{ILP wall} (extracting more
instruction-level parallelism costs too much hardware for small gains) and the
\textbf{memory wall} (DRAM did not keep up with cores). Multiprocessors give higher
throughput within power limits and some availability (one processor can fail);
their limits are Amdahl's law, contention for shared resources and tasks that cannot
be parallelized.

\begin{eg}[Amdahl's law limits parallel speedup]
To reach a speedup of 80 with 100 processors: $80=\dfrac{1}{f/100+(1-f)}$ gives
$f=0.9975$: only 0.25\% of the original computation may be
sequential.\HP{\S5.1, Example}{349} Group~4: $f=0.90$ on 64 processors gives
$1/(0.1+0.9/64)\approx8.8$; $f=0.99$ gives $\approx39.3$. Removing dependences
raises $f$.
\end{eg}

\begin{eg}[The cost of communication]
32 processors, remote memory access 200 ns, clock 3.3 GHz (0.3 ns), base $\CPI=0.5$
with all references local; 0.2\% of instructions make a remote reference costing
$200/0.3\approx666$ cycles. $\CPI=0.5+0.2\%\times666\approx1.7$: the all-local
machine is $1.7/0.5=3.4\times$ faster.\HP{\S5.1, Example}{350}
\end{eg}

\subsubsection{Memory organization}
\begin{center}
\small
\begin{tabular}{@{}P{3.6cm}P{5.2cm}P{5.2cm}@{}}
\toprule
Organization & Structure & Implication \\
\midrule
Centralized shared memory: \textbf{SMP}, \textbf{UMA} & all processors share one
memory (and I/O) with the same latency & simple sharing; memory bandwidth limits
scaling (about 8--32 cores) \\
Distributed shared memory: \textbf{DSM}, \textbf{NUMA} & memory distributed among
nodes, one shared address space; local accesses faster than remote & scales; remote
latency, directory coherence \\
Message passing (clusters) & private memories per node & explicit messages \\
\bottomrule
\end{tabular}
\end{center}
In an SMP (symmetric multiprocessor) all processors are equals and can run OS and
application code; multicore chips are SMPs. UMA examples: a multicore SoC (Apple
Silicon); NUMA: multi-socket database servers.\HP[348]{\S5.1, Figures~5.1--5.2}{346}

\subsubsection{Flynn's classification}
Classify by the number of \emph{instruction streams} and \emph{data streams}:
\begin{center}
\small
\begin{tabular}{@{}lP{11cm}@{}}
\toprule
SISD & single instruction, single data: a uniprocessor (it may still be pipelined or
superscalar) \\
SIMD & single instruction, multiple data: one instruction on many data items; data-level
parallelism (vector units, multimedia extensions, GPUs) \\
MISD & multiple instruction, single data: \textbf{no commercial multiprocessor of
this type has been built} \\
MIMD & multiple instruction, multiple data: each processor fetches its own
instructions and data; task-level parallelism; the most general and common \\
\bottomrule
\end{tabular}
\end{center}
\HP{\S1.2}{10}
\begin{remark}
\slidefix{} The Part~1 slides give pipelining as an application of MISD. A
pipeline processes one instruction stream; H\&P states that no commercial MISD
machine exists. The number of cores alone does not determine the Flynn class.
\end{remark}

% =============================================================================
\subsection{Synchronization}
\label{sec:sync}

\begin{definition}[Race condition, critical section]
A \textbf{race condition}: the result depends on the uncontrolled order in which
threads access shared data, with at least one write and no synchronization
(\texttt{count = count + 1} is a read, an add and a write; two threads can lose an
update). The \textbf{critical section} is the code that accesses shared data; a
correct protocol guarantees \textbf{mutual exclusion} (at most one thread inside),
\textbf{progress} (if it is free and threads want in, one enters) and
\textbf{bounded waiting} (no starvation).
\end{definition}

\paragraph{Atomic primitives.} A lock built from ordinary loads and stores has its
own race (two threads both see the flag free). Hardware provides an atomic
read-modify-write:\HP[388]{\S5.5}{387}
\begin{itemize}
  \item \textbf{Atomic exchange} (swap): swap a register and a memory word;
  \textbf{test-and-set}: return the old value and set it to 1;
  \textbf{fetch-and-increment}; \textbf{compare-and-swap}.
  \item \textbf{Load linked / store conditional} (\texttt{LL}/\texttt{SC}): the
  store succeeds (returns 1) only if the location was not written since the
  \texttt{LL}; the pair builds any of the above (Section~\ref{sec:mips}).
\end{itemize}

\begin{definition}[Spin lock]
A lock that a processor tries to acquire in a loop until it succeeds. Used when the
lock is held very briefly. Simple version (\texttt{R1} $=$ lock address, 0 $=$ free):
\begin{lstlisting}
        DADDUI R2,R0,#1
lockit: EXCH   R2,0(R1)    ; atomic exchange
        BNEZ   R2,lockit   ; already locked?
\end{lstlisting}
Release: store 0. Each exchange is a write, so with coherent caches every attempt
invalidates the others' copies and floods the bus. The \textbf{test-and-test-and-set}
version spins on a cached \emph{read} and tries the exchange only when the lock
looks free:\HP[390]{\S5.5}{389}
\begin{lstlisting}
lockit: LD     R2,0(R1)    ; load of lock
        BNEZ   R2,lockit   ; not available: spin
        DADDUI R2,R0,#1    ; load locked value
        EXCH   R2,0(R1)    ; swap
        BNEZ   R2,lockit   ; branch if lock wasn't 0
\end{lstlisting}
\end{definition}

\begin{center}
\small
\begin{tabular}{@{}lP{5.6cm}P{5.6cm}@{}}
\toprule
& Spin lock & Mutex (sleeping lock) \\
\midrule
While waiting & loops, CPU busy & sleeps; the OS runs other work \\
Acquire cost & immediate & context switch \\
Best for & very short waits (microseconds) & long waits (disk, network) \\
One core & useless (the owner cannot run) & works \\
\bottomrule
\end{tabular}
\end{center}
\begin{itemize}
  \item \textbf{Deadlock}: threads wait for each other in a circle (needs mutual
  exclusion, hold and wait, no preemption, circular wait); avoid it e.g.\ by taking
  locks in one global order.
  \item \textbf{At large scale}: test-and-test-and-set, \textbf{exponential backoff}
  (wait 1, 2, 4, \dots{} before retrying), \textbf{queue locks} (MCS: each waiter
  spins on its own flag; fair; Linux qspinlock).
\end{itemize}

\begin{definition}[Barrier]
A \textbf{barrier} forces all threads to wait until every thread has arrived (e.g.\
between program phases). Simple version: a lock-protected counter; the last arrival
resets the counter and sets \texttt{release}; the others spin on \texttt{release}.
Used in a loop it can deadlock: a fast thread reaches the barrier again and clears
\texttt{release} before a slow thread has left. A \textbf{sense-reversing barrier}
fixes this: each thread flips a private \texttt{local\_sense} and waits for
\texttt{release == local\_sense}.\HPA[15]{I}{\S I.3, Figures~I.2--I.3}{14} Large
machines use \textbf{tree barriers} to avoid one hot counter.
\end{definition}

% =============================================================================
\subsection{Memory Consistency}
\label{sec:consistency}

\begin{definition}[Coherence vs.\ consistency]
\textbf{Coherence} defines what values a read of \emph{one} location can return.
\textbf{Consistency} defines \emph{when} a written value must be seen by reads,
i.e.\ the allowed order of accesses to \emph{different} locations. Coherent caches
alone do not fix that order.\HP{\S5.2}{353}
\end{definition}

\begin{definition}[Sequential consistency (SC)]
The result of any execution is the same as if the operations of all processors were
executed in some sequential order, and the operations of each processor appear in
program order (Lamport). All four orderings W$\to$R, R$\to$R, R$\to$W, W$\to$W are
kept.\HP[395]{\S5.6}{392}
\end{definition}

\begin{eg}[Store buffering]
Initially $x=y=0$. Thread 0: \texttt{x = 1; r0 = y}. Thread 1: \texttt{y = 1; r1 = x}.
Under SC, $r0=r1=0$ is impossible (whichever store is first in the global order is
seen by the other thread's load). With a store buffer (relaxing W$\to$R), each store
can wait in its buffer while the following load reads the old value, so $r0=r1=0$
can happen. Coherence does not prevent it: it is about two different locations.
\end{eg}

\paragraph{Relaxed models.} Relaxing orderings lets hardware hide write latency;
programs then use synchronization (fences, acquire/release) to restore order.
\begin{center}
\small
\begin{tabular}{@{}lP{5.8cm}P{4.4cm}@{}}
\toprule
Model & Relaxes & Used in \\
\midrule
TSO / processor consistency & W$\to$R: a load may pass an earlier store (write buffer); stores stay in order & x86 \\
PSO (partial store order) & W$\to$R and W$\to$W (stores to different addresses) & SPARC (PSO mode) \\
Weak ordering & all orders between synchronization operations; syncs act as fences & ARM, PowerPC style \\
Release consistency & as weak ordering, but distinguishes \emph{acquire} (later ops wait for it) and \emph{release} (earlier ops finish before it) & lock-based programs \\
\bottomrule
\end{tabular}
\end{center}

% =============================================================================
\subsection{Dependences and Instruction-Level Parallelism}
\label{sec:ilp}

\subsubsection{Dependences}
\begin{definition}[Kinds of dependence]
\begin{itemize}
  \item \textbf{Data (true) dependence}: $j$ uses a result of $i$ (directly or through
  a chain). The only real dependence; hazard \textbf{RAW}. Cannot be removed.
  \item \textbf{Name dependence}: two instructions use the same register or memory
  location without data flowing: \textbf{antidependence} ($j$ writes what $i$
  reads, hazard \textbf{WAR}) and \textbf{output dependence} (both write, hazard
  \textbf{WAW}). Removed by \textbf{renaming}.
  \item \textbf{Control dependence}: whether $j$ executes depends on a branch;
  handled by prediction and speculation.
  \item \textbf{Resource dependence}: two operations need the same unit (structural).
\end{itemize}
A dependence is a property of the program; whether it becomes a hazard (a stall)
depends on the pipeline.\HP[155]{\S3.1}{150}
\end{definition}

\begin{eg}[Renaming (Group 4)]
\texttt{I1: R1 = R2 + R3}; \texttt{I2: R4 = R1 * 2}; \texttt{I3: R2 = R5 - 1};
\texttt{I4: R4 = R6 + R7}. I1$\to$I2 is true (RAW on R1); I1$\to$I3 is anti (WAR on
R2); I2$\to$I4 is output (WAW on R4). Renaming I3's R2 to P2 and I4's R4 to P4
removes the last two; I1$\to$I2 remains.
\end{eg}

\paragraph{Finding parallelism (program level).}
\begin{itemize}
  \item \textbf{Bernstein's conditions} (Hwang; \beyond{} H\&P): with input set $I_i$
  and output set $O_i$, $P_1\parallel P_2$ iff $I_1\cap O_2=\emptyset$ (no anti),
  $I_2\cap O_1=\emptyset$ (no flow) and $O_1\cap O_2=\emptyset$ (no output).
  $\parallel$ is not transitive.
  \item A \textbf{dependence graph} has statements as nodes; its longest path (the
  \emph{critical path}) is the minimum time with unlimited processors.
  \item A \textbf{loop-carried dependence}: iteration $i$ uses a value from iteration
  $i-d$. No carried dependence $\Rightarrow$ the loop is parallel (DOALL); carried but
  no cycle $\Rightarrow$ restructure; a cycle $\Rightarrow$ DOACROSS
  (post/wait synchronization) or serial.
  \item \textbf{GCD test}: for a store to $X[a\,i+b]$ and a load from $X[c\,i+d]$, a
  loop-carried dependence requires $\gcd(c,a)$ to divide $d-b$. Example:
  \texttt{X[2*i+3] = X[2*i] * 5.0}: $\gcd(2,2)=2$ does not divide $-3$, so no
  dependence.\HP{\S4.5}{319}
  \item \textbf{Reductions} (a sum) can be regrouped as a tree: $N-1$ dependent adds
  become $\log_2N$ steps.
\end{itemize}

\begin{eg}[Bernstein's conditions (Group 4, after Hwang)]
P1: $C=D\times E$; P2: $M=G+C$; P3: $A=B+C$; P4: $C=L+M$; P5: $F=G\div E$. Only
P1$\parallel$P5, P2$\parallel$P3, P2$\parallel$P5, P3$\parallel$P5 and P4$\parallel$P5
hold. P1$\parallel$P5 and P5$\parallel$P2, but P1 and P2 are not parallel (flow on $C$).
\end{eg}

\subsubsection{Compiler scheduling and loop unrolling}
To avoid a stall, a dependent instruction must be separated from its producer by
the producer's latency. With latencies load$\to$FP ALU 1, FP ALU$\to$FP ALU 3, FP
ALU$\to$store 2 and one stall between an integer ALU op and a branch, the loop
\texttt{x[i] = x[i] + s} takes 9 cycles per element unscheduled and 7 scheduled.
\textbf{Loop unrolling} replicates the body (with different registers), removes
repeated loop overhead and gives the scheduler independent instructions: unrolled 4
times and scheduled, it takes 14 cycles for 4 elements (3.5 per element), at the
cost of more registers and code.\HP[161]{\S3.2, Figure~3.2}{157}

\subsubsection{Superpipelining and superscalar processors}
\begin{itemize}
  \item \textbf{Superpipelining}: more, shorter stages, so a faster clock; still at
  most one instruction issued per cycle; larger branch and load penalties (e.g.\ the
  8-stage MIPS R4000).\HPA{C}{\S C.6}{61}
  \item \textbf{Superscalar}: issue several instructions per cycle (IPC $>1$); the
  hardware checks dependences; \emph{statically} scheduled (in order) or
  \emph{dynamically} scheduled (out of order). Dependences, branches, resources and
  issue width limit the gain.
  \item Both can be combined.
\end{itemize}

\begin{eg}[Superpipeline (Part 1 exercise)]
A 12-stage pipeline issues at most one instruction per cycle; an instruction takes
2 ns unpipelined. It is scalar but superpipelined. For 100 instructions (balanced
stages, no stalls): unpipelined $100\times2=200$ ns; 5 stages
$(5+99)\times0.4=41.6$ ns; 12 stages $(12+99)\times2/12=18.5$ ns.
\end{eg}

\subsubsection{Issue and completion policies}
\begin{center}
\small
\begin{tabular}{@{}P{3.8cm}P{4.8cm}P{5.4cm}@{}}
\toprule
Policy & Behaviour & Hazards to check, cost \\
\midrule
In-order issue, in-order completion & a stalled instruction blocks all later ones & RAW only; simple, slow with long operations \\
In-order issue, out-of-order completion & short operations may finish before long ones & RAW and WAW; a reorder buffer (ROB) restores order for precise exceptions \\
Out-of-order issue, out-of-order completion & any ready instruction in a window may start & RAW, WAR, WAW; needs renaming, an instruction window, a ROB; fastest, most expensive \\
\bottomrule
\end{tabular}
\end{center}
\textbf{Dynamic scheduling} (Tomasulo's algorithm) issues in order but executes out
of order, renaming registers through reservation stations to remove WAR and WAW.
\textbf{Hardware speculation} adds a reorder buffer: results \emph{commit} in
order, so exceptions are precise and mispredicted work can be
discarded.\HP[175]{\S3.4}{167}\HP{\S3.6}{184}

\subsubsection{VLIW}
\begin{definition}[Very long instruction word]
A VLIW processor issues a fixed number of operations packed by the
\textbf{compiler} into one long instruction (one slot per functional unit). It has no
hardware hazard detection or dynamic scheduling, so the hardware is simpler and uses
less power. The compiler must find the parallelism: \textbf{loop unrolling},
\textbf{software pipelining} (each new iteration takes instructions from different
original iterations), \textbf{trace scheduling} and \textbf{superblocks} (scheduling
across branches along the likely path).\HP[196]{\S3.7}{193}
\end{definition}
Drawbacks: unused slots hold \texttt{nop}s, so code grows; latencies are fixed in the
code, so binaries are not compatible across implementations; a stall in one unit
stalls the whole bundle; poor fit for unpredictable branches and pointer-heavy code.

\begin{center}
\small
\begin{tabular}{@{}P{3.3cm}P{2cm}P{2.4cm}P{3.2cm}P{2.6cm}@{}}
\toprule
Approach & Issue & Hazard detection & Scheduling & Examples \\
\midrule
Superscalar (static) & dynamic & hardware & static & embedded MIPS, ARM \\
Superscalar (dynamic) & dynamic & hardware & dynamic & few today \\
Superscalar (speculative) & dynamic & hardware & dynamic, speculation & Intel Core i7, AMD \\
VLIW/LIW & static & mainly software & static & TI C6x \\
EPIC & mainly static & mainly software & mostly static & Itanium \\
\bottomrule
\end{tabular}
\end{center}
\HP{\S3.7, Figure~3.15}{194}

% =============================================================================
\subsection{Multithreading}
\label{sec:multithreading}

\begin{definition}[Hardware multithreading]
Several threads share one core's functional units; the core keeps a separate PC,
register set and page table per thread, so it can switch threads without an OS
context switch. It raises \emph{throughput} and utilization by hiding stalls; a single
thread is not faster.\HP[226]{\S3.12}{223}
\end{definition}
\begin{itemize}
  \item \textbf{Coarse-grained}: switch only on costly stalls (L2/L3 misses). Little
  slowdown of each thread, but the pipeline must refill on a switch, so short stalls
  are not hidden.
  \item \textbf{Fine-grained}: switch threads every cycle (round robin), skipping
  stalled threads. Hides short and long stalls, but slows each individual thread.
  \item \textbf{Simultaneous multithreading (SMT)}: on a multiple-issue, dynamically
  scheduled core, issue instructions from \emph{several threads in the same cycle}.
  It exploits \textbf{thread-level and instruction-level parallelism} together and
  fills the issue slots one thread leaves empty (horizontal waste); Intel
  Hyper-Threading.
\end{itemize}

\begin{eg}[Hiding a miss (Part 1)]
A single-issue core runs thread T1, which misses in the cache for 3 cycles. Alone: T1,
T1 (miss), idle, idle, idle, T1: 3 instructions in 6 cycles. With a second thread the
idle cycles run T2: 6 instructions in 6 cycles. T1 still finishes at cycle 6.
\end{eg}

% =============================================================================
\subsection{Cache Coherence}
\label{sec:coherence}

\begin{eg}[The coherence problem]
$X=1$ in memory; write-through caches. A reads $X$ (1); B reads $X$ (1); A stores 0
into $X$: A's cache and memory hold 0, but B's cache still holds 1, and B will read
the stale value.\HP{\S5.2, Figure~5.3}{352}
\end{eg}

\begin{definition}[Coherence]
A memory system is \textbf{coherent} if (1) a read by P of $X$ after P wrote $X$, with
no other writes in between, returns P's value (\emph{program order}); (2) a read of
$X$ sufficiently after another processor wrote it returns the written value (a
\emph{coherent view}); (3) writes to the same location are \textbf{serialized}: all
processors see them in the same order.\HP[353]{\S5.2}{352}
\end{definition}
Coherent caches allow \textbf{migration} (move data to the local cache: less latency
and traffic) and \textbf{replication} (copies in several caches for simultaneous
reads), done in hardware because software is too slow.

\begin{center}
\small
\begin{tabular}{@{}lP{6cm}P{6cm}@{}}
\toprule
& Snooping & Directory-based \\
\midrule
Sharing state & kept by every cache holding a copy & kept in one place per block (a directory, distributed with memory) \\
Communication & broadcast on a bus; all controllers snoop & point-to-point messages to the home node \\
Write serialization & the bus orders writes & the directory orders them \\
Scales to & small SMPs (a few dozen cores) & hundreds of processors (DSM/NUMA) \\
Cost & bus bandwidth, broadcast energy & directory storage, 3-hop misses \\
\bottomrule
\end{tabular}
\end{center}
\HP[355]{\S5.2}{354}

\paragraph{Write invalidate vs.\ write update.} \textbf{Write invalidate} (used in
practice): a writer first gets exclusive access by invalidating all other copies;
later writes are local. \textbf{Write update} (broadcast): every write sends the new
value to all copies, which uses much more bandwidth.\HP[356]{\S5.2}{355}

\subsubsection{MSI and MESI}
\begin{definition}[MSI protocol]
A write-invalidate protocol with three states per cache block:
\textbf{M}odified (only copy, dirty; may read and write), \textbf{S}hared (clean,
read-only, possibly in other caches), \textbf{I}nvalid.\HP[361]{\S5.2, Figures~5.5--5.7}{358}
\end{definition}
\begin{center}
\small
\begin{tabular}{@{}lll@{}}
\toprule
Event & Transition & Bus action \\
\midrule
Local read miss & I $\to$ S & read miss (BusRd) \\
Local write (miss or hit in S) & I/S $\to$ M & write miss / invalidate \\
Remote read of a block in M & M $\to$ S & write back (flush), supply data \\
Remote write (or invalidate) & S, M $\to$ I & M writes back first \\
Replacement of an M block & M $\to$ I & write back \\
\bottomrule
\end{tabular}
\end{center}

\begin{definition}[MESI protocol]
MESI adds \textbf{E}xclusive: the block is in only this cache and is clean. A read
miss with no other copy loads the block in E. A later write changes E $\to$ M
\textbf{without any bus transaction}, which optimizes a block read and then written
by the same core (e.g.\ single-threaded code). A remote read changes E $\to$ S. Intel
i7 uses MESIF (adds Forward); MOESI adds Owned.\HP{\S5.2}{362}
\end{definition}

\begin{eg}[MSI trace (Group 4)]
Two processors, memory $X=0$, write-back caches.
\begin{center}
\small
\begin{tabular}{@{}llccc@{}}
\toprule
Step & Bus & Cache A & Cache B & Memory \\
\midrule
A reads $X$ & read miss & S 0 & I & 0 \\
B reads $X$ & read miss & S 0 & S 0 & 0 \\
A writes 1 & invalidate & M 1 & I & 0 \\
B reads $X$ & miss; A writes back & S 1 & S 1 & 1 \\
A writes 2 & invalidate & M 2 & I & 1 \\
B writes 0 & write miss; A writes back, invalidates & I & M 0 & 2 \\
\bottomrule
\end{tabular}
\end{center}
\end{eg}
\begin{remark}
\slidefix{} Group~4's table leaves memory at 1 after ``B writes 0''. In H\&P's
protocol a block in M that sees a write miss on the bus is written back before it
is invalidated, so memory holds 2; B's new value 0 stays in B's cache until it is
written back.\HP{\S5.2, Figure~5.6}{360}
\end{remark}

\begin{eg}[MESI trace (Group 4)]
Three processors, $X=0$. A reads: miss, no other copy, A in \textbf{E}. A writes 1:
\textbf{silent} E $\to$ M. B reads: A writes back, A and B in S (memory 1). C reads:
A, B, C in S. C writes 2: invalidate, A and B in I, C in M. Under MSI, A's first
write would have needed an invalidate on the bus.
\end{eg}

\subsubsection{Directory protocols}
Each memory block has a directory entry with a state and a \textbf{sharing bit
vector} (one bit per node): \textbf{Uncached} (no copy), \textbf{Shared} (one or more
clean copies; memory up to date), \textbf{Modified} (exactly one node, the owner, has
written it; memory stale). Messages go point to point: local node (requester), home
node (memory and directory), remote node (has a copy). A write to a shared block
makes the directory send invalidates only to the nodes in the bit
vector.\HP[383]{\S5.4, Figure~5.20}{379}

\paragraph{The fourth C: coherence misses.} \textbf{True sharing}: a word written by
one processor is read by another. \textbf{False sharing}: processors use
\emph{different} words in the same block, so invalidations move the block with no
real data flow; fix with padding or a different data layout.\HP{\S5.3}{366}

% =============================================================================
\subsection{Worked Problems}
\label{sec:mp-problems}

\begin{eg}[Flynn (Part 1)]
One core, one instruction at a time: SISD. A vector unit applying one operation to
8 elements: SIMD. Four processors running different branches of a simulation on
different data: MIMD. ``Four cores'' alone: not enough information.
\end{eg}

\begin{eg}[Which primitive? (Part 1)]
Workers updating one shared counter: a lock (spinning wastes CPU while waiting).
No worker may start phase 2 before all finish phase 1: a barrier (the slowest worker
delays all). A lock held 50~\textmu s with 16 runnable threads on 4 cores: spinning
burns about $12\times50=600$ thread-\textmu s; sleeping is better. 1024 workers with
a barrier after every tiny step: barrier contention dominates; use a tree barrier or
coarser phases. 8 threads add to a total, then one prints it: a lock around
\texttt{total += x} and a barrier before printing.
\end{eg}

\begin{eg}[2-way superscalar issue (Part 1)]
\texttt{I1: R1 = R2 + R3}; \texttt{I2: R4 = R5 + R6}; \texttt{I3: R7 = R1 + R8};
\texttt{I4: R9 = R10 + R11}. I3 needs I1's result. In-order 2-way issue:
(I1, I2), then (I3, I4): 2 issue cycles. A width of 2 does not guarantee 2 per cycle:
if I2 depended on I1, 3 cycles would be needed.
\end{eg}

% =============================================================================
\subsection{Textbook Exercises}
\label{sec:mp-tbx}

\begin{remark}
Group~4 worked H\&P Exercises 3.14 and 5.4 in their slides. Check two points before
reusing those answers: (1) H\&P counts a latency of $n$ as $n$ stall cycles between
producer and consumer (\S3.2, p.~158), while the slides insert $n-1$; (2) Exercise
5.4 starts from the cache state of Figure~5.35 (P0 already holds block 110 in M,
P3 holds 120 in S), while the slides assume empty caches.
\end{remark}

\begin{tbexercise}{5.1}{412}{\S5.2 (snooping MSI), \hppages{}{356}{361}}
Each of the four-block, direct-mapped, two-word-block private caches is kept
coherent with the MSI snooping protocol. Initial state (state, tag, data):
\begin{center}
\scriptsize
\begin{tabular}{@{}lccc@{}}
\toprule
& P0 & P1 & P3 \\
\midrule
B0 & I 100 (00 10) & I 100 (00 10) & S 120 (00 20) \\
B1 & S 108 (00 08) & M 128 (00 68) & S 108 (00 08) \\
B2 & M 110 (00 30) & I 110 (00 10) & I 110 (00 10) \\
B3 & I 118 (00 10) & S 118 (00 18) & I 118 (00 10) \\
\bottomrule
\end{tabular}
\end{center}
Memory: 100: 00 10, 108: 00 08, 110: 00 10, 118: 00 18, 120: 00 20, 128: 00 28,
130: 00 30. Applying each action independently to this state, give the blocks that
change and the value returned by each read: (a) P0: read 120; (b) P0: write 120
$\leftarrow$ 80; (c) P3: write 120 $\leftarrow$ 80; (d) P1: read 110; (e) P0: write
108 $\leftarrow$ 48; (f) P0: write 130 $\leftarrow$ 78; (g) P3: write 130 $\leftarrow$ 78.
\end{tbexercise}

\begin{tbexercise}{5.4}{415}{\S5.2--5.3, \hppages{}{356}{367}}
Same initial state as Exercise 5.1. A miss costs $N_{memory}=100$ cycles if served by
memory and $N_{cache}$ if served by another cache (40 in implementation 1, 130 in
implementation 2); a write hit that invalidates costs $N_{invalidate}=15$; a
write-back adds $N_{writeback}=10$. Compute the total stall cycles for the MSI
protocol (the book also asks for MOSI): (a) P0: read 110, P3: read 110, P0: read
110; (b) P1: read 120, P3: read 120, P0: read 120; (c) P0: write 120 $\leftarrow$ 80,
P3: read 120, P0: read 120; (d) P0: write 108 $\leftarrow$ 88, P3: read 108, P0: write
108 $\leftarrow$ 98.
\end{tbexercise}

\begin{tbexercise}{5.7}{416}{\S5.5 (spin locks), \hppages{}{389}{391}}
P0, P1 and P3 try to acquire a lock at address 0x100, with the cache state of
Exercise 5.1, implementation-1 timing and 1000-cycle critical sections. (a) With the
simple exchange spin lock, about how many memory stall cycles does each processor
incur before acquiring the lock? (b) Repeat with the test-and-test-and-set lock.
(c), (d) About how many interconnect transactions occur with each lock?
\end{tbexercise}

\begin{tbexercise}{5.8}{417}{\S5.6 (consistency), \hppages{}{392}{395}}
Under SC a processor must stall a read until an earlier write is ordered; TSO lets
reads pass the processor's own buffered writes. One memory operation per cycle; hits
and write-buffer hits cost nothing; misses cost as in Exercise 5.4 (implementation 1).
With the cache state of Exercise 5.1, how many stall cycles occur before each
operation under SC and under TSO? (a) P0: write 110 $\leftarrow$ 80, P0: read 108;
(b) P0: write 100 $\leftarrow$ 80, P0: read 108; (c) P0: write 110 $\leftarrow$ 80, P0:
write 100 $\leftarrow$ 90; (d) P0: write 100 $\leftarrow$ 80, P0: write 110
$\leftarrow$ 90.
\end{tbexercise}

\begin{tbexercise}{5.19}{426}{\S5.1 and \S1.9 (Amdahl), \hppages{}{348}{350}}
$F(i,p)$ is the fraction of time that exactly $i$ of $p$ processors are usable
($\sum_{i=1}^{p}F(i,p)=1$), and with $i$ processors the application runs $i$ times
faster. Rewrite Amdahl's law to give the speedup as a function of $p$.
\end{tbexercise}

\begin{tbexercise}{5.29}{428}{\S5.5, \hppages{}{387}{388}}
Implement the classical test-and-set instruction using the load-linked/store-conditional
pair.
\end{tbexercise}

\begin{tbexercise}{3.14}{255}{\S3.2 and \S3.7, \hppages{}{157}{162},
\hppages{}{193}{197}}
The DAXPY loop $Y=aX+Y$ (100 elements; \texttt{R1}, \texttt{R2} point to $X$, $Y$):
\begin{lstlisting}
      DADDIU R4,R1,#800
foo:  L.D    F2,0(R1)
      MUL.D  F4,F2,F0
      L.D    F6,0(R2)
      ADD.D  F6,F4,F6
      S.D    F6,0(R2)
      DADDIU R1,R1,#8
      DADDIU R2,R2,#8
      DSLTU  R3,R1,R4
      BNEZ   R3,foo
\end{lstlisting}
Latencies: FP multiply $\to$ FP ALU 6; FP add $\to$ FP ALU 4; FP multiply $\to$
store 5; FP add $\to$ store 4; integer operations and loads $\to$ any 2. One-cycle
delayed branch resolved in ID; full bypassing. (a) Single issue: show the loop
unscheduled and scheduled, the cycles per element of $Y$, and how much faster the
clock would have to be to match the scheduling gain. (b) Unroll the loop as many
times as needed to remove all stalls; show the schedule and the cycles per element.
(The book's part (c) schedules it on a 5-slot VLIW.)
\end{tbexercise}
